 \documentclass[times,10pt]{article}
%\documentclass{sig-alternate}
\input{globaldefs}
%%%%%%%%%%%%%%%%%%%%%%%%%%%%%%%%%%%%%%%%%%%%%%%%%%%%%%%%%%%%%%%%%%%%%%%%%%%%%
%   Begin the document.                             %
%%%%%%%%%%%%%%%%%%%%%%%%%%%%%%%%%%%%%%%%%%%%%%%%%%%%%%%%%%%%%%%%%%%%%%%%%%%%%
\begin{document}

\newcommand{\gasnet}[0]{GASNet\ }





%%%%%%%%%%%%%%%%%%%%%%%%%%%%%%%%%%%%%%%%%%%%%%%%%%%%%%%%%%%%%%%%%%%%%%%%%%%%%
%   Title                                   %
\title{\gasnet Teams Specification and Design Document}
\author{Rajesh Nishtala, Dan Bonachea, Paul Hargrove\\
    \emph{rajeshn@cs.berkeley.edu, bonachea@cs.berkeley.edu, PHHargrove@lbl.gov} \\\\
    Computer Science Division, University of California at Berkeley \\ 
    Future Technologies Group, Lawrence Berkeley National Lab\\ \\
    Preliminary Work: For Internal Distribution Only}
\maketitle
%%%%%%%%%%%%%%%%%%%%%%%%%%%%%%%%%%%%%%%%%%%%%%%%%%%%%%%%%%%%%%%%%%%%%%%%%%%%%


\section{Introduction}
In a variety of applications the use of collective communication is very
important for both expressing the communication pattern of the application as
well as achieving the best performance. However, one of the fundamental aspects
that makes the collectives useful is the ability to use them on a subset (or
\textit{team}) of the processors rather than all the processors in the program.
This document will focus on the design and implementation of these teams within
\gasnet as well as the API that exposes them to applications or libraries that
use \gasnet. 

Since \gasnet allows hierarchical declaration of nodes and images we must be
careful about the terminology we use and how we compare our implementation with
MPI's. Henceforth, we will use ``node" to refer to a \gasnet process. Each of
these \gasnet nodes can potentially contain many  \gasnet images (which can be
implemented on top of different threads if one chooses). The distinction allows
one to perform shared memory communication optimizations across various images
within the same node since they will share the same address space. In MPI
parlance this hierarchy is not explicit and thus all the threads of control
visibile to the application are at the same level which we will call ``tasks."
Our teams will consist of \gasnet images rather than nodes since we want
different threads in the same process to be a part of different teams. We will
allow images to have their own relative ranks within different teams but the
\gasnet node identifier will remain fixed throughout the life of the program. 

\section{Design Goals}
\begin{itemize}
\item One of the important and fundamental characteristics of \gasnet that
needs to be addressed in the design of teams is the ability to have multiple
``images" of \gasnet on the same \gasnet process, thus allowing one to perform
shared memory/intra-node optimizations for collectives whenever possible. This
ability will need to also be effectively exposed in the design of the teams. 

\item Each team of images will need to be sufficiently isolated from the others
so that the communication and synchronization mechanisms will never interfere
with those from another team. This will allow libraries to be written on top of
\gasnet that take teams as arguments and be guaranteed that operations on one
team will be isolated from operations performed on another team provided that
they do not share members. If they do, then special care must be done at the
application level.

\item The team creation and destruction must be abstract and expressive enough
to easily and quickly build up whatever subset of processors the user wishes.

\item The teams must fit seamlessly into the existing \gasnet collectives
framework so that collective performance on all the images in a \gasnet program
is not hindered by the addition of teams. 

\end{itemize}

\section{Similarities and Differences with MPI}
The problem of communication with a subset of the processors is certainly not a
new one and has been around in the MPI specification for many years. In MPI
parlance the notion of teams is expressed as communicators. Many of the ideas
presented in the API for MPI communicators (and groups) are orthogonal to their
use of two-sided communication and their design of the collectives and thus a
lot of these ideas are applicable here. However, there are some important
differences that will need to be addressed as well. 

\subsection{Similarities}
\begin{itemize}
\item \textbf{Communicators:} As stated above, the notion of an MPI
communicator is very similar to our concept of teams, modulo some of the
differences stated below. A team, like a communicator, is a set of \gasnet
images along with a preallocated set of buffer space and communication meta
data (e.g. a distributed buffer space manager) to allow fast collective
communication across the various members of the team. 

\item \textbf{Groups:} We will also adopt the idea of an MPI Group into
\gasnet. The primary function of a \gasnet group is to allow easy team
construction. The process of creating a team is an expensive process due to the
setup of all the required meta-data. Since a group merely an ordered set of
images without the additional meta data necessary for communication and
synchronization their construction and modification can be done using much
simpler, and therefore more efficient, operations. Thus the general model, like
MPI, will be to build up a group based on the different operations provided in
the API and then construct a team around a group once the group has been
finalized. 
  
%\item \textbf{API:} Our initial API for both teams and groups have been
%heavily adopted from the MPI. However as we will show there are a few
%differences that need to be addressed. 

\end{itemize}

\subsection{Differences}
\begin{itemize}
\item \textbf{\gasnet Images:} As stated above the biggest difference
(especially in the implementation) that needs to be addressed is the difference
between a node an image and how these relate to the teams. MPI does not have
this problem since there is no notion of a hierarchical structure between MPI
tasks; they are all at the same level. 

\item \textbf{Scratch Space:} \gasnet has an explicit segment which allows one
to take advantage of some very important communication optimizations (e.g.
RDMA) in one-sided operations. Since we wish to leverage some of these same
optimizations in the collectives, the collectives (and thus teams) will need to
reclaim some of the space that has been allocated to the end user to provide
the best possible performance. The auxiliary scratch space for the initial
GASNET\_TEAM\_ALL will be handled directly within \gasnet so  that it isn't
visible to the end user, however further team construction will necessitate the
user explicitly managing these buffers.

%Note that since \gasnet is indented as a compiler target, this will probably
%not be exposed to final end users and will only be relevant to writers of a
%runtime

\item \textbf{Usage:} Another important and distinct difference is that MPI
allows one to use a communicator for point-to-point messaging such as sends and
receives. However, we will disallow applications that use \gasnet from using
teams for point-to-point communication. The reason for this decision is that it
breaks the up the roles of teams strictly for use in collectives rather than
the full generality of point-to-point communication. Thus each image will have
one globally unique name for point-to-point communication. The translation
routines that are provided so that one can specify the root for rooted
collectives relative to the other images in the team. 

\end{itemize}


\newcommand{\group}[0]{gasnet\_group\_t}
\newcommand{\grouparg}[0]{gasnet\_group\_t\ g}


\newcommand{\team}[0]{gasnet\_team\_t}
\newcommand{\teamarg}[0]{gasnet\_team\_t\ t}

\newcommand{\node}[0]{gasnet\_node\_t}
\newcommand{\nodearg}[0]{gasnet\_node\_t\ n}

\newcommand{\image}[0]{gasnet\_image\_t}
\newcommand{\imagearg}[0]{gasnet\_image\_t\ i}

\section{Groups}
A \gasnet group is defined as an ordered set of \gasnet images that will take
part in the construction of a \gasnet team. A group is designed to be simple
and easy to construct based on other groups. All these operations are
non-collective. The following is a list of the operations that will operate on
and work with a \group. 

\begin{itemize}
\item Constant \texttt{\image\ GASNET\_IMAGE\_UNDEFINED 0xffffff} \\ 
A constant defined for an invalid image number. 

\item GlobalVar: \texttt{\group\ gasnet\_empty\_group} \\ 
A global variable reserved for an empty group to aid in group creation. 

\item \texttt{\image\ gasnet\_group\_size(\grouparg)}\\ 
Returns the total number of images in a \texttt{\group} $g$. Note that
\texttt{\image} is a wrapper around an integer that is wide as the maximum
count of the total \gasnet images allowed.  

\item \texttt{\image\ gasnet\_group\_size\_on\_node(\grouparg, \nodearg)} \\ 
Returns the number of images on node $n$ that are a part $g$. 

\item \texttt{\image\ gasnet\_group\_my\_image(\grouparg)} \\ 
Returns my relative image rank on $g$. 

\item \texttt{void gasnet\_group\_translate\_images(\group\ A, \group\ B,\\ 
              \image\ *input\_vec, \image\ *output\_vec, \image\ vec\_len)} \\ 
Given a set of input images in $input\_vec$ from group $A$ the function will
translate them to their corresponding relative ranks in $B$. If some images do
not exist in node $B$, then they the array slot will be filled with
GASNET\_IMAGE\_UNDEFINED.

\item \texttt{\node\ gasnet\_group\_image\_to\_node(\grouparg, \imagearg)} \\ 
Given a relative rank $i$ in $g$, the function returns the node $n$ that
contains the image. 

\item \texttt{void gasnet\_group\_node\_to\_images(\grouparg, \\ 
              \nodearg, \image\ *outimages)} \\ 
Fills the output array $outimages$ with the images of $g$ that are on $n$. Use
\texttt{gasnet\_group\_size\_on\_node} to get the count of how many elements
this function will return. 

\item \texttt{enum\{GASNET\_GROUP\_IDENT , GASNET\_GROUP\_SIMILAR, GASNET\_GROUP\_UNEQUAL\} \\ 
           gasnet\_group\_compare(\group\ a, \group\ b)} \\ 
This function returns one of three values:  
\begin{enumerate}
\item If $a$ and $b$ have the same size and ordering of images then the
function will return $GASNET\_GROUP\_IDENT$. 

\item If $a$ and $b$ contain the same number and set of images but the relative
ranks of the images are different then the function returns
$GASNET\_GROUP\_SIMILAR$.

\item Otherwise the function will return $GASNET\_GROUP\_UNEQUAL$.

\end{enumerate}

\item \texttt{\group\ gasnet\_group\_from\_team(\teamarg)} \\ Returns a
\texttt{\group} that represents the images in $t$. 

\item \texttt{\group\ gasnet\_group\_incl(\grouparg, \image\ * in\_images, \image\ image\_count)} \\ 
Returns a group that contains all the images in $g$ along with the images in
$in\_images$ appended to the end. 

\item \texttt{\group\ gasnet\_group\_excl(\grouparg, \image\ * in\_images, \image\ image\_count)} \\ 
Returns a \texttt{\group} that removes the images contained in $in\_images$
from $g$. The ordering of the remaining images from $g$ will not change. 

\item \texttt{\group\ gasnet\_group\_diff(\group\ A, \group\ B)} \\ 
Returns a \texttt{\group} that is the set difference of $A$ and $B$. The
resultant set will be all the images of $A$ that are not in $B$, ordered by
their relative rankings in $A$. 

\item \texttt{\group\ gasnet\_group\_intersect(\group\ A, \group\ B)} \\ 
Returns a \texttt{\group} that is the set intersection of $A$ and $B$. The
images are ordered according to their ordering in $A$. 

\item \texttt{\group\ gasnet\_group\_union(\group\ A, \group\ B)} \\ 
Returns a \texttt{\group} that is the set union of $A$ and $B$.  The images of
$A$ are placed before the images of $B$ in the ordering. 

\item \texttt{size\_t gasnet\_group\_to\_node\_seg\_size(\grouparg)} \\ 
Given a group $g$ the function returns the amount of scratch space a team
constructed with $g$ will need on the node that the function is called on. 

\item \texttt{size\_t gasnet\_split\_to\_node\_seg\_size(\teamarg, \image\ color, \image\ relrank)} \\ 
This is the only group-like function that must be a collective call. It will
return the amount of space required for scratch on the calling node in order to
perform a split of the team $t$ with the corresponding $color$ and $relrank$
arguments that one would use for the \texttt{gasnet\_team\_split()} function
that will be relrank below.  If there are multiple teams that will be created
on the same node, the function will return a valid size for the scratch space
for one representative image from each of the teams that will reside on the
node. A nonzero return value indicates that this image has been chosen to
allocate scratch space for the corresponding team split. 

%\item GroupRangeIncl (maybe: scalability) ???
%\item GroupRangeExcl (maybe: scalability) ???
\item \texttt{void gasnet\_group\_free(\grouparg)} \\ 
This function will perform the required housekeeping to free all the state
associated with $g$.  

\end{itemize}



\section{Team Operators}
Team Creation is collective across the parent group.
\begin{itemize}

\item GlobalVar: \texttt{GASNET\_TEAM\_ALL}  \\ 
Initial team that all images are a part of. Comparable to MPI\_COMM\_WORLD. 

\item \texttt{\team\ gasnet\_team\_split(\teamarg, \\ 
              \image\ color, \image\ relrank, void *local\_seg\_space)} \\ 
This function splits the team $t$ based on the $color$ and $relrank$ arguments.
All images that call the function with the same $color$ argument will be
separated into the same team. The $relrank$ argument specifies the relative
position of the image within the new team. If multiple images call the split
function with the same color and relrank argument the results are undefined.
Note that an image is responsible for allocating the local\_seg\_space if the
corresponding call to \texttt{gasnet\_split\_to\_node\_seg\_size()} returned a
nonzero value. One can pass GASNET\_IMAGE\_UNDEFINED to the $relrank$ to
indicate that the image is not to be part of any team. 

\item \texttt{\team\ gasnet\_team\_from\_group(\teamarg, \\\grouparg, void\ *local\_seg\_space)} \\ 
This function will create a child team of $t$ from the group $g$ with the given
scratch space.  The results are undefined if all members of $g$ are not part of
$t$. 

%\item (TeamSplit is sufficient for Team Duplicate)

\item \texttt{\image\ gasnet\_team\_my\_image(\teamarg)} \\ 
Given a team the function returns a relative rank of the calling image within
$t$. If the calling image is not containted in $t$ then
GASNET\_IMAGE\_UNDEFINED is returned. 

\item \texttt{\image\ gasnet\_team\_size(\teamarg)} \\ 
Returns the number of images in $t$. 

\item \texttt{void* gasnet\_team\_free(\teamarg)} \\ 
Free a team by cleaning up all associated resources with the team. Once the
free routine is called it is safe to free the scratch space that was passed in
to create the team. As a convenience, the original argument scratch space
argument is returned so that a free can be called on the returned value. 

%\item No TeamCompare (can just do pointer comparison on Teams)
\end{itemize}




\section{Notes to Self About Implementation.}
\begin{enumerate}
\item Keep Team Hierarchy
\item Reference Count Groups
\end{enumerate}

\end{document}
